\begin{enumerate}[label=\textbf{\ansnum.\arabic*}]
  \item $6\cdot1+2\cdot2+60\cdot1=70$ bytes. Without \texttt{<}, \texttt{struct} uses native alignment: the two \texttt{H} must start at an even offset; after $6$ bytes they already do, so here the size stays $70$. For \verb!"4s4B3ff2f"! it also stays $32$. But a layout like \verb!"B H"! would become $4$ bytes instead of $3$; always write \texttt{<} for network packets so the result does not depend on the computer.
  \item
\begin{verbatim}
def our_team(state, team_number):
    for k in (0, 1):
        if state.teams[k]["number"] == team_number:
            return k
    raise ValueError(f"team {team_number} is not in this game")
\end{verbatim}
  \item Index $2$: player numbers start at $1$, list indices at $0$.
  \item \verb!encode_return(3, 12, False, (-2500.0, 1000.0, 0.3), 0.4, (1200.0, -300.0))! (player $3$ and team $12$ as examples).
  \item $256/2=128$\,s. Compute the difference modulo $256$: \verb!lost = (new - old - 1) % 256!; with packets arriving in order this is the number of missing packets even across the wrap from $255$ to $0$.
\end{enumerate}
